\section{Institutional and Empirical Bridge}
\label{sec:institutional}

The model abstracts from any single national program. Its two instruments correspond to a recurring distinction in place-based green industrial policy between firm-facing support for decarbonizing investment and area-level investment in productive infrastructure such as clean-power access, grid capacity, hydrogen or utility networks, ports, industrial land, and related common inputs. The theoretical mapping is functional rather than literal: $s_i$ changes the private return to green investment, while $h_i$ changes local production conditions and can also affect territorial emissions.

Japan's recent GX industrial-location policy illustrates the coexistence of these margins. The Ministry of Economy, Trade and Industry's GX Strategic Area framework organizes regional industrial-cluster development around assets such as existing industrial complexes and geographically concentrated decarbonized power, and explicitly includes integrated power and communications infrastructure \citep{METIGX2026}. The same framework separately contains a firm-facing ``decarbonized-power regional contribution'' investment-support channel. The corresponding 2026 program subsidizes large electricity-user investment conditional on the use of decarbonized electricity and contribution to host power-source communities, and states industrial-competitiveness strengthening as an objective \citep{METIGXSubsidy2026}.

These programs support the paper's distinction between place-based productive infrastructure and firm-facing green investment support, but they should not be read as evidence that Japanese local governments independently choose both instruments exactly as in the theoretical game. The model is a stylized interjurisdictional-policy game. Japan provides institutional motivation for the relevant margins, not a literal mapping of legal authority or financing incidence.

More generally, the OECD reports that subnational governments have an increasing implementation role in Japan's environmental strategy and that about 60 percent of Japanese local governments have committed to net zero by 2050 \citep{OECDJapan2025}. This supports the assumption that territorial decarbonization objectives can matter at the subnational level, while leaving the paper agnostic about the precise delegation mechanism behind $\bar E_i$ and $d$.

The parameter $\mu>0$ is likewise a scope condition rather than a claim that all environmental spending is privately profitable. The mechanism is most relevant for green investment that also affects product-market competitiveness---for example through energy efficiency, process modernization, reduced energy-cost exposure, or product attributes valued by customers. If green investment is pure compliance expenditure with no competitiveness effect, Corollary~\ref{cor:mu-zero} shows that the cross-jurisdictional product-market transmission from a rival green subsidy disappears.

\subsection{Empirical predictions}

The theory generates testable implications, but they are threshold and local-comparative-static predictions rather than a globally linear response model.

First, the sign of the local infrastructure response to rival firm-facing green support can differ across product-substitutability regimes. In the canonical model the response is negative below an interior threshold and positive above it, but it is not monotonically increasing throughout the low-rivalry region. An empirical design should therefore allow nonlinearities or threshold interactions rather than impose a globally positive coefficient on ``rival subsidy $\times$ competition.''

Second, in the canonical neighborhood the switching threshold is lower where shared green infrastructure has a larger productivity effect. Exposure to electricity-intensive production, hydrogen or steam networks, grid bottlenecks, port access, or other common industrial inputs can provide empirical variation related to $\nu$.

Third, again locally around the canonical region, switching occurs at a lower rivalry level when ordinary productivity investment is easier to adjust, corresponding to a lower $k_x$. Measures of capital adjustment, automation intensity, or the responsiveness of conventional investment to demand conditions could therefore be used to characterize this margin.

Fourth, the response should weaken sharply where environmental investment has little effect on marginal cost or product-market position, the empirical counterpart of $\mu\simeq0$.

Finally, on the interior reduced system the rival policies enter jurisdiction $A$ through the composite
\[
y_B=\frac{\mu}{k_g}s_B+\nu h_B.
\]
Hence a reduced-form policy-reaction coefficient by itself does not identify a subsidy-specific structural channel. Empirical work should distinguish instrument-specific shocks---for example through eligibility rules, program discontinuities, or staggered changes---and should separate predetermined exposure to rival jurisdictions from endogenous policy matching.

These predictions point toward a spatial policy-reaction design rather than a simple regression of emissions on subsidies. The present paper remains theoretical; the empirical discussion is intended to clarify which observed reactions would be diagnostic of the model and which would also arise under generic interjurisdictional competition.